% ============================================================================
% DM-2026-013: Marking a Software Product Form that binds a product to a payload
% ============================================================================
\newcommand{\UniqueID}{DM-2026-013}
\newcommand{\DocumentDate}{September 26, 2026}
\newcommand{\AuthorName}{Bruce Dombrowski}
\newcommand{\AuthorTitle}{Systems Engineer, Payload Software Integration}
\newcommand{\ToField}{PSI Team, Payload Developers, Submitting Organizations, NASA JSC CUI Program Manager}
\newcommand{\SubjectField}{Marking Guidance: Export Classification on a completed Software Product Form}
\newcommand{\OPRField}{Payload Software Integration (PSI)}

\newcommand{\dmContent}{%
\begin{enumerate}

\dmsection{Purpose}

\item This memorandum gives a payload developer a rule for filling the
\textbf{Export Classification} and \textbf{Export Statement} fields on a
Software Product Form (SPF), and says who adjudicates when the answer is not
obvious. It is prompted by a concrete case: how a developer should mark an SPF
for a secure shell (SSH) implementation delivered with a payload.

\dmsection{Background}

\item SSP 57061, Paragraph 6.1.4, requires a Test SPF and a Flight SPF for a
payload software update, alongside a Version Description Document and malware
scan evidence. SPFs are entered in the ORBIT Software Products (SWP)
application.

\item The SPF carries an Export Classification field and an Export Statement
field. The field definitions note that the Payload Software Integration
Engineer (PSIE) may enter them. Neither the field definitions nor the
application states \emph{when} a product raises a classification, or what a
submitter should weigh. Both fields default to an unset value, and an unset
value is indistinguishable from a considered answer of ``none''.

\item Security-relevant products are precisely the products most likely to
appear on an SPF and most likely to matter here. SSH is one: SSP 57000
Appendix I and SSP 57003 Appendix L each give SSH Key Delivery and SSH
Configuration their own requirement clauses, alongside Encrypted Protocols.
Antivirus software and operating systems are others.

\dmsection{Analysis}

\item \textbf{Neither fact is sensitive alone.} A commercial software version
is published by its vendor. The existence of an ISS payload is public. A
submitter reasoning about either in isolation will correctly conclude that
nothing needs marking.

\item \textbf{The form binds them, and the binding is a different fact.} A
completed SPF states that a \emph{named} on-orbit system is running a
\emph{specific} version of a security-relevant product. That is a statement
about which published weaknesses that system carries, which is not derivable
from either half.

\item \textbf{The programme has already taken this view once.} DM-2026-008
treats antivirus distribution specifics as controlled for the same reason: the
sensitivity is created by associating security software with the systems it
protects, not by the software's identity. A payload-bound SSH version is the
same shape of fact.

\item \textbf{This does not disturb DM-2026-007 or DM-2026-012.} Those
determinations concern blank forms and field definitions --- material that
shows \emph{which fields exist} without populating them. That reasoning is
unaffected and is reaffirmed here. What changes at the point of population is
the aggregation, not the form.

\item \textbf{Over-marking has a cost too.} If every SPF is marked, marking
stops carrying information and the fields become noise. The rule below is
therefore scoped to security-relevant products rather than applied to all.

\dmsection{Decision}

\item A submitter \textbf{shall not leave Export Classification unset} on a
submitted SPF. An unset value is not a determination.

\item Where the product is \textbf{security-relevant} --- including but not
limited to secure shell implementations, antivirus software, operating systems,
and cryptographic libraries --- the submitter shall \textbf{refer the marking to
the PSIE} rather than self-determine, and shall treat the completed form as
potentially controlled until the PSIE responds.

\item Where the product is \textbf{not security-relevant}, the submitter may
enter the classification directly, and should record the basis in the Export
Statement field rather than leaving it empty.

\item Questions that the PSIE cannot resolve shall be referred to the NASA JSC
CUI Program Manager, following the route DM-2026-008 established.

\dmsection{Implementation}

\item Payload developer training material shall state that Export
Classification is a required judgement rather than an optional field, and shall
use SSH as the worked example, since it is the case where the Software Product
Form and the IT security requirements meet.

\item PSI recommends that the SWP field definitions carry a short paragraph
under Export Classification stating that the classification of a form is not
the classification of the product, because the form binds the product to a
payload. That recommendation is raised separately against the next revision of
those definitions.

\dmsection{Conditions that change this guidance}

\item A determination by the NASA JSC CUI Program Manager that payload-bound
software version information is, or is not, controlled --- which would replace
the referral rule above with a direct answer.

\item A change to the SWP application that removes the unset default, or that
adds its own marking guidance.

\item Any change to DM-2026-008's treatment of security software associated
with the systems it protects.

\dmsection{References}

\item SSP 57061, Standard Payload Integration Agreement for ISS Payloads,
Paragraph 6.1.4, Software Updates.

\item SSP 57000 Appendix I and SSP 57003 Appendix L, Information Technology
System Security.

\item ORBIT SWP field definitions, Software Product Form.

\item DM-2026-007, blank form CUI determination; DM-2026-008, desk instruction
CUI classification; DM-2026-012, training package CUI determination.

\end{enumerate}
}

% ============================================================================
% Decision Memorandum Base Template
% ============================================================================
% This is the shared base template for all Decision Memoranda.
% DO NOT edit this file for individual decisions.
%
% Usage: Create a thin wrapper that defines variables and content, then:
%   % ============================================================================
% Decision Memorandum Base Template
% ============================================================================
% This is the shared base template for all Decision Memoranda.
% DO NOT edit this file for individual decisions.
%
% Usage: Create a thin wrapper that defines variables and content, then:
%   % ============================================================================
% Decision Memorandum Base Template
% ============================================================================
% This is the shared base template for all Decision Memoranda.
% DO NOT edit this file for individual decisions.
%
% Usage: Create a thin wrapper that defines variables and content, then:
%   \input{_template.tex}
%
% Required variables (define before \input):
%   \UniqueID       - e.g., DM-2026-001
%   \DocumentDate   - e.g., January 19, 2026
%   \AuthorName     - Signer name
%   \AuthorTitle    - Signer title
%   \ToField        - Distribution
%   \SubjectField   - Subject line
%   \OPRField       - Office of Primary Responsibility
%   \dmContent      - The full content (enumerate environment with \dmsection items)
%
% ============================================================================

\documentclass[11pt,letterpaper]{article}

% ============================================================================
% PACKAGES
% ============================================================================
\usepackage[margin=1in, top=1.15in, headheight=30pt]{geometry}
\usepackage{graphicx}
\usepackage{fancyhdr}
\usepackage{lastpage}
\usepackage{enumitem}
\usepackage{xcolor}
\usepackage{hyperref}
\usepackage{tabularx}

% ============================================================================
% DOCUMENT CONFIGURATION
% ============================================================================
\hypersetup{
    colorlinks=true,
    linkcolor=blue,
    urlcolor=blue,
    pdfborder={0 0 0}
}

% Remove paragraph indentation
\setlength{\parindent}{0pt}
\setlength{\parskip}{6pt}

% List formatting - match Word style
\setlist[enumerate]{leftmargin=0.5in, labelwidth=0.25in, labelsep=0.1in, align=left, itemsep=12pt, topsep=12pt}
\setlist[enumerate,1]{label=\arabic*.}
\setlist[itemize]{leftmargin=0.5in, itemsep=3pt, topsep=6pt}

% ============================================================================
% HEADER AND FOOTER CONFIGURATION
% ============================================================================
\pagestyle{fancy}
\fancyhf{}

% Header: text-only, Boeing / ISS Program identification
\fancyhead[L]{\large\textbf{Decision Memorandum}\\[-2pt]\scriptsize Boeing \textbar{} International Space Station Program}
\fancyhead[R]{\small \UniqueID\\[-2pt]\scriptsize \DocumentDate}
\renewcommand{\headrulewidth}{0.4pt}

% Footer: Unique ID (left), Page X of Y (center), Date (right)
\fancyfoot[L]{\small \UniqueID}
\fancyfoot[C]{\small Page \thepage\ of \pageref{LastPage}}
\fancyfoot[R]{\small \DocumentDate}
\renewcommand{\footrulewidth}{0.4pt}

% ============================================================================
% CUSTOM COMMANDS
% ============================================================================
% Bold section headers for numbered sections
\newcommand{\dmsection}[1]{\item \textbf{#1}\par}

% ============================================================================
% BEGIN DOCUMENT
% ============================================================================
\begin{document}

% ----------------------------------------------------------------------------
% UNIQUE ID (top right)
% ----------------------------------------------------------------------------
\hfill \UniqueID

\vspace{0.25in}

% ----------------------------------------------------------------------------
% SIGNATURE BLOCK
% ----------------------------------------------------------------------------
\underline{\hspace{2.5in}}\\[2pt]
\AuthorName\\
\AuthorTitle

\vspace{0.25in}

% ----------------------------------------------------------------------------
% MEMO HEADER FIELDS
% ----------------------------------------------------------------------------
\textbf{DATE:} \DocumentDate

\textbf{TO:} \ToField

\textbf{SUBJECT:} \SubjectField

\textbf{OFFICE OF PRIMARY RESPONSIBILITY:} \OPRField

\vspace{0.25in}

% ============================================================================
% MAIN CONTENT - NUMBERED SECTIONS
% ============================================================================
\dmContent

\end{document}

%
% Required variables (define before \input):
%   \UniqueID       - e.g., DM-2026-001
%   \DocumentDate   - e.g., January 19, 2026
%   \AuthorName     - Signer name
%   \AuthorTitle    - Signer title
%   \ToField        - Distribution
%   \SubjectField   - Subject line
%   \OPRField       - Office of Primary Responsibility
%   \dmContent      - The full content (enumerate environment with \dmsection items)
%
% ============================================================================

\documentclass[11pt,letterpaper]{article}

% ============================================================================
% PACKAGES
% ============================================================================
\usepackage[margin=1in, top=1.15in, headheight=30pt]{geometry}
\usepackage{graphicx}
\usepackage{fancyhdr}
\usepackage{lastpage}
\usepackage{enumitem}
\usepackage{xcolor}
\usepackage{hyperref}
\usepackage{tabularx}

% ============================================================================
% DOCUMENT CONFIGURATION
% ============================================================================
\hypersetup{
    colorlinks=true,
    linkcolor=blue,
    urlcolor=blue,
    pdfborder={0 0 0}
}

% Remove paragraph indentation
\setlength{\parindent}{0pt}
\setlength{\parskip}{6pt}

% List formatting - match Word style
\setlist[enumerate]{leftmargin=0.5in, labelwidth=0.25in, labelsep=0.1in, align=left, itemsep=12pt, topsep=12pt}
\setlist[enumerate,1]{label=\arabic*.}
\setlist[itemize]{leftmargin=0.5in, itemsep=3pt, topsep=6pt}

% ============================================================================
% HEADER AND FOOTER CONFIGURATION
% ============================================================================
\pagestyle{fancy}
\fancyhf{}

% Header: text-only, Boeing / ISS Program identification
\fancyhead[L]{\large\textbf{Decision Memorandum}\\[-2pt]\scriptsize Boeing \textbar{} International Space Station Program}
\fancyhead[R]{\small \UniqueID\\[-2pt]\scriptsize \DocumentDate}
\renewcommand{\headrulewidth}{0.4pt}

% Footer: Unique ID (left), Page X of Y (center), Date (right)
\fancyfoot[L]{\small \UniqueID}
\fancyfoot[C]{\small Page \thepage\ of \pageref{LastPage}}
\fancyfoot[R]{\small \DocumentDate}
\renewcommand{\footrulewidth}{0.4pt}

% ============================================================================
% CUSTOM COMMANDS
% ============================================================================
% Bold section headers for numbered sections
\newcommand{\dmsection}[1]{\item \textbf{#1}\par}

% ============================================================================
% BEGIN DOCUMENT
% ============================================================================
\begin{document}

% ----------------------------------------------------------------------------
% UNIQUE ID (top right)
% ----------------------------------------------------------------------------
\hfill \UniqueID

\vspace{0.25in}

% ----------------------------------------------------------------------------
% SIGNATURE BLOCK
% ----------------------------------------------------------------------------
\underline{\hspace{2.5in}}\\[2pt]
\AuthorName\\
\AuthorTitle

\vspace{0.25in}

% ----------------------------------------------------------------------------
% MEMO HEADER FIELDS
% ----------------------------------------------------------------------------
\textbf{DATE:} \DocumentDate

\textbf{TO:} \ToField

\textbf{SUBJECT:} \SubjectField

\textbf{OFFICE OF PRIMARY RESPONSIBILITY:} \OPRField

\vspace{0.25in}

% ============================================================================
% MAIN CONTENT - NUMBERED SECTIONS
% ============================================================================
\dmContent

\end{document}

%
% Required variables (define before \input):
%   \UniqueID       - e.g., DM-2026-001
%   \DocumentDate   - e.g., January 19, 2026
%   \AuthorName     - Signer name
%   \AuthorTitle    - Signer title
%   \ToField        - Distribution
%   \SubjectField   - Subject line
%   \OPRField       - Office of Primary Responsibility
%   \dmContent      - The full content (enumerate environment with \dmsection items)
%
% ============================================================================

\documentclass[11pt,letterpaper]{article}

% ============================================================================
% PACKAGES
% ============================================================================
\usepackage[margin=1in, top=1.15in, headheight=30pt]{geometry}
\usepackage{graphicx}
\usepackage{fancyhdr}
\usepackage{lastpage}
\usepackage{enumitem}
\usepackage{xcolor}
\usepackage{hyperref}
\usepackage{tabularx}

% ============================================================================
% DOCUMENT CONFIGURATION
% ============================================================================
\hypersetup{
    colorlinks=true,
    linkcolor=blue,
    urlcolor=blue,
    pdfborder={0 0 0}
}

% Remove paragraph indentation
\setlength{\parindent}{0pt}
\setlength{\parskip}{6pt}

% List formatting - match Word style
\setlist[enumerate]{leftmargin=0.5in, labelwidth=0.25in, labelsep=0.1in, align=left, itemsep=12pt, topsep=12pt}
\setlist[enumerate,1]{label=\arabic*.}
\setlist[itemize]{leftmargin=0.5in, itemsep=3pt, topsep=6pt}

% ============================================================================
% HEADER AND FOOTER CONFIGURATION
% ============================================================================
\pagestyle{fancy}
\fancyhf{}

% Header: text-only, Boeing / ISS Program identification
\fancyhead[L]{\large\textbf{Decision Memorandum}\\[-2pt]\scriptsize Boeing \textbar{} International Space Station Program}
\fancyhead[R]{\small \UniqueID\\[-2pt]\scriptsize \DocumentDate}
\renewcommand{\headrulewidth}{0.4pt}

% Footer: Unique ID (left), Page X of Y (center), Date (right)
\fancyfoot[L]{\small \UniqueID}
\fancyfoot[C]{\small Page \thepage\ of \pageref{LastPage}}
\fancyfoot[R]{\small \DocumentDate}
\renewcommand{\footrulewidth}{0.4pt}

% ============================================================================
% CUSTOM COMMANDS
% ============================================================================
% Bold section headers for numbered sections
\newcommand{\dmsection}[1]{\item \textbf{#1}\par}

% ============================================================================
% BEGIN DOCUMENT
% ============================================================================
\begin{document}

% ----------------------------------------------------------------------------
% UNIQUE ID (top right)
% ----------------------------------------------------------------------------
\hfill \UniqueID

\vspace{0.25in}

% ----------------------------------------------------------------------------
% SIGNATURE BLOCK
% ----------------------------------------------------------------------------
\underline{\hspace{2.5in}}\\[2pt]
\AuthorName\\
\AuthorTitle

\vspace{0.25in}

% ----------------------------------------------------------------------------
% MEMO HEADER FIELDS
% ----------------------------------------------------------------------------
\textbf{DATE:} \DocumentDate

\textbf{TO:} \ToField

\textbf{SUBJECT:} \SubjectField

\textbf{OFFICE OF PRIMARY RESPONSIBILITY:} \OPRField

\vspace{0.25in}

% ============================================================================
% MAIN CONTENT - NUMBERED SECTIONS
% ============================================================================
\dmContent

\end{document}

